\section{Involución de hojas en \(4\) representaciones}

La involuci\'on \(C^*\mapsto-C^*\) organiza cuatro respuestas distintas:
\begin{equation}
\label{eq:modular-sheet-parities}
\begin{aligned}
S_{\rm bi}(A,-C^*)&=S_{\rm bi}(A,C^*),\\
\Gamma_{6\to8}(A,-C^*)&=-\Gamma_{6\to8}(A,C^*),\\
(\widehat\mu,\widehat\varepsilon)&\longleftrightarrow
(\widehat\varepsilon,\widehat\mu),\\
\widehat c&\longmapsto\widehat c,
\qquad D_A\longmapsto D_A.
\end{aligned}
\end{equation}
La informaci\'on bicapa es par; la orientaci\'on areal es impar; las dos
constitutivas forman un doblete, y los determinantes quedan fijos. Esta tabla
de paridades precede a cualquier interpretaci\'on como simetr\'ia f\'isica.

\section{Transformaci\'on racional involutiva en el espacio de par\'ametros CKM}

Sea
\begin{equation}
\label{eq:modular-Mckm}
M_{\rm CKM}=
\begin{pmatrix}
2&-5&28\\
0&7&-25/2\\
1&-20&-2/3
\end{pmatrix},
\qquad
\det M_{\rm CKM}=-\frac{3857}{6}.
\end{equation}
La inversa existe sobre \(\mathbb Q\) y es
\begin{equation}
\label{eq:modular-Mckm-inverse}
M_{\rm CKM}^{-1}=
\begin{pmatrix}
1528/3857&3380/3857&801/3857\\
75/3857&176/3857&-150/3857\\
6/551&-30/551&-12/551
\end{pmatrix}.
\end{equation}
Con
\begin{equation}
\label{eq:modular-ckm-coordinates}
h=\frac{C^*}{6},
\qquad
\Delta=\frac{180}{\pi}\Delta_4,
\qquad
\boldsymbol\theta=M_{\rm CKM}(A,h,\Delta)^T,
\end{equation}
la coordenada impar entra en la direcci\'on
\begin{equation}
\label{eq:modular-ch}
c_h=\begin{pmatrix}-5\\7\\-20\end{pmatrix}.
\end{equation}

Si \(c_A=(2,0,1)^T\) y
\(c_\Delta=(28,-25/2,-2/3)^T\), las tres columnas
\((c_A,c_h,c_\Delta)\) son precisamente la base transportada por
\(M_{\rm CKM}\). La segunda fila de
\cref{eq:modular-Mckm-inverse} es el covector \(\ell_h\) que extrae la
coordenada impar:
\begin{equation}
\label{eq:modular-dual-coordinate-check}
\ell_hc_A=0,
\qquad
\ell_hc_h=1,
\qquad
\ell_hc_\Delta=0.
\end{equation}

\begin{theorem}[Reflexi\'on CKM inducida por la hoja]
\label{thm:modular-ckm-reflection}
Sea
\begin{equation}
\label{eq:modular-Dh-lh}
D_h=\operatorname{diag}(1,-1,1),
\qquad
\ell_h=\frac1{3857}\begin{pmatrix}75&176&-150\end{pmatrix}.
\end{equation}
La involuci\'on inducida en el espacio de \(\boldsymbol\theta\) es
\begin{equation}
\label{eq:modular-Rckm}
\boxed{
\mathcal R_{\rm CKM}
=M_{\rm CKM}D_hM_{\rm CKM}^{-1}
=I-2c_h\ell_h.}
\end{equation}
Satisface
\begin{equation}
\label{eq:modular-Rckm-properties}
\mathcal R_{\rm CKM}^2=I,
\qquad
\det\mathcal R_{\rm CKM}=-1,
\qquad
\mathcal R_{\rm CKM}M_{\rm CKM}=M_{\rm CKM}D_h.
\end{equation}
\end{theorem}

\begin{proof}
La multiplicaci\'on de
\cref{eq:modular-Mckm,eq:modular-Mckm-inverse} da la identidad; en
particular, \cref{eq:modular-dual-coordinate-check} prueba que
\(c_h\ell_h\) es el proyector de rango uno sobre \(\mathbb Qc_h\) a lo
largo de \(\operatorname{span}\{c_A,c_\Delta\}\). El producto exacto
\(\ell_hc_h=1\) implica
\[
(I-2c_h\ell_h)^2
=I-4c_h\ell_h+4c_h(\ell_hc_h)\ell_h=I.
\]
El operador fija punto por punto el hiperplano
\(\ker\ell_h=\operatorname{span}\{c_A,c_\Delta\}\) y cambia el signo de
\(c_h\). Sus autovalores son, por tanto, \(1,1,-1\); de aqu\'i
\(\det\mathcal R_{\rm CKM}=-1\) y
\(\Tr\mathcal R_{\rm CKM}=1\). Finalmente,
\[
\mathcal R_{\rm CKM}M_{\rm CKM}
=M_{\rm CKM}D_hM_{\rm CKM}^{-1}M_{\rm CKM}
=M_{\rm CKM}D_h,
\]
que prueba la identidad de entrelazamiento.
\end{proof}

La matriz expl\'icita es
\begin{equation}
\label{eq:modular-Rckm-matrix}
\mathcal R_{\rm CKM}=
\begin{pmatrix}
4607/3857&1760/3857&-1500/3857\\
-150/551&199/551&300/551\\
3000/3857&7040/3857&-2143/3857
\end{pmatrix}.
\end{equation}
Aunque \(\mathcal R_{\rm CKM}\) no es una reflexi\'on ortogonal para el
producto eucl\'ideo de las tres coordenadas angulares, s\'i lo es para la
m\'etrica racional transportada
\begin{equation}
\label{eq:modular-ckm-metric}
G_{\rm CKM}=(M_{\rm CKM}^{-1})^TM_{\rm CKM}^{-1},
\qquad
\mathcal R_{\rm CKM}^TG_{\rm CKM}\mathcal R_{\rm CKM}=G_{\rm CKM}.
\end{equation}
En efecto, en la base \((c_A,c_h,c_\Delta)\), la m\'etrica es la identidad
y la involuci\'on es \(D_h\). Por eso el t\'ermino ``reflexi\'on'' tiene un
contenido m\'etrico preciso y no describe s\'olo que
\(\mathcal R_{\rm CKM}^2=I\).
La fase \(\delta_{\rm CP}=9A\) permanece fija bajo esta reflexi\'on. Por
ello, \(\mathcal R_{\rm CKM}\) no es la conjugaci\'on CP f\'isica, sino el
transporte de la inversi\'on de hoja HMT al espacio de mezcla.

La inversa racional de \(M_{\rm CKM}\) recupera
\begin{align}
A&=\frac{1528\theta_{12}+3380\theta_{23}+801\theta_{13}}{3857},
\label{eq:modular-A-from-ckm}\\
C^*&=\frac{6(75\theta_{12}+176\theta_{23}-150\theta_{13})}{3857}.
\label{eq:modular-C-from-ckm}
\end{align}
La tercera coordenada tambi\'en se recupera sin p\'erdida:
\begin{equation}
\label{eq:modular-Delta-from-ckm}
\Delta
=\frac{6\theta_{12}-30\theta_{23}-12\theta_{13}}{551}.
\end{equation}
Estas ecuaciones son cambios de coordenadas. No convierten CKM en causa
retrospectiva de \(A,C^*\) o de \(\gammaH\).

Como reconocimiento terminal, el constructor previo determina
\begin{equation}
\label{eq:modular-ckm-values}
(\theta_{12},\theta_{23},\theta_{13},\delta_{\rm CP})
=(13.003313589509^\circ,
2.398056157332^\circ,
0.213977382244^\circ,
65.676173123554^\circ).
\end{equation}
La reflexi\'on racional se demuestra antes de usar estos decimales.

\begin{proposition}[Resultado principal]
La transici\'on \(6\to8\) se formaliza primero mediante la inclusi\'on
\(\iota_{6,8}:\mathcal F_6\hookrightarrow\mathcal F_8\), de la que se
obtienen
\(90,120\) y el defecto \(-1/12\). La palabra de seis trits proporciona
despu\'es la biyecci\'on \(729\leftrightarrow729\), y
\(\mathcal J_{6\to8}\) transporta los dos modos colectivos al borde. Sobre
ese antecedente, el \`area y el Hamiltoniano modular reconstruyen exactamente
las dos familias; la purificaci\'on termocampo doble determina la entrop\'ia
de entrelazamiento, y la primera ley
finita conserva el defecto de entrop\'ia relativa. En CKM, la inversi\'on de
hoja se convierte en una reflexi\'on racional de rango impar uno.
\end{proposition}

\begin{remark}[Control negativo]
La genealog\'ia inicial falla si \(\iota_{6,8}\) no es inyectiva, si los
conteos no son \(90/120\), si \(\delta_{6,8}^{\rm inc}\neq-1/12\), si
\(\beta_{729}\) no es biyectiva o si se la presenta falsamente como
isomorfismo aditivo. El transporte colectivo falla si
\(\mathcal J_{6\to8}D_{\rm inc}\neq
D_{\rm area}\mathcal J_{6\to8}\). La reconstrucci\'on falla si el
determinante de los observables de borde no es treinta o si
\cref{eq:modular-inverse-channels} no recupera las ocupaciones.
La TFD falla si su traza parcial no es \(\rho_A\). La primera ley no se
extiende a cambios finitos sin la entrop\'ia relativa de
\cref{eq:modular-finite-first-law}. El tipo
\(III_{\lambda_\partial}\) queda invalidado si un canal no persiste, si no hay
producto compatible o si cambia el grupo de razones. La reflexi\'on CKM
queda refutada por cualquiera de
\(\mathcal R^2\neq I\), \(\det\mathcal R\neq-1\) o
\(\mathcal RM\neq MD_h\). Una f\'ormula
\(\boldsymbol\theta=f(\gammaH)\) sin entrelazador viola la independencia
tipada de las dos realizaciones funcionales del antecedente com\'un.
\end{remark}

\subsection*{Alcance lógico y dominio de validez}
Inclusi\'on \(H\subset O\), conteos \(90/120\), dos normalizaciones del
defecto \(-1/12\), biyecci\'on de seis trits \(729\leftrightarrow729\),
entrelazador colectivo, escala \(a_0\), operadores finitos, espectro areal,
TFD, defecto orientado, primera ley local y finita y reflexi\'on CKM:
\emph{Teorema interno exacto con sus estructuras de partida expl\'icitas}.
Clasificaci\'on
\(III_{\lambda_\partial}\): \emph{Teorema con hip\'otesis expl\'icitas de
persistencia}. Interpretaci\'on de la copia termocampo doble como exterior f\'isico:
\emph{Realizaci\'on f\'isica condicionada}. La magnitud \(\gammaH\) se adopta
como \emph{Resultado antecedente demostrado}; no se rederiva en este volumen.

